\section{Challenges and Future Directions}
\label{sec:challenges}

Despite rapid progress, current systems remain far from general Physical AI. 
We summarize several open challenges and future directions.

\paragraph{From language priors to dense physical grounding.}
LLM-based world knowledge is broad but sparse. 
Future systems must ground language priors into dense representations of geometry, motion, contact, force, uncertainty, and temporal dynamics.

\paragraph{From VLMs to physically faithful perception.}
VLMs can describe visual scenes, but Physical AI requires physically faithful perception that supports spatial reasoning, affordance understanding, temporal prediction, and action feasibility.

\paragraph{From VLAs to generalist embodied policies.}
VLA models connect perception and language to action, but their generalization is limited by embodiment-specific action spaces and scarce robot data. 
Future work should study scalable action representations, cross-embodiment transfer, and policy learning from heterogeneous data.

\paragraph{From video generation to world modeling.}
Not every video generation model is a world model. 
Physical AI requires models that are action-conditioned, temporally consistent, controllable, and physically plausible. 
Bridging photorealistic generation with physical correctness remains a major challenge.

\paragraph{From simulation to real-world deployment.}
World models and simulators can support training and planning, but sim-to-real transfer remains difficult. 
Future systems must address distribution shift, sensor noise, embodiment mismatch, safety constraints, and closed-loop robustness.

\paragraph{From closed frontier systems to reproducible science.}
Many frontier Physical AI systems are closed-source or only partially documented. 
This creates challenges for fair comparison, reproducibility, and scientific progress. 
Developing transparent evaluation protocols for closed systems is an important direction for the community.